
To install Orbiter, we first check if git is installed. Type
$$
\verb'git'
$$
If you get an error message, please install git. 
If you get a long output with commands, you are fine. 
Git is installed. 

\bigskip


The steps to install Orbiter differ slightly, 
depending on your operating system.
For Windows, we need ``Windows Subsystem Linux.''
On Linux/unix based systems, including Macintosh, 
you will utilize a terminal window. 
You will also need the development tools. 
You need the C++ compiler and library tools.
On Macintosh, you need Xcode and the command-line tools. 
You can get them for free on the App-store.
You will also need the compiler 
generator tools flex and bison.
They are open source projects. On Macintosh, you can use the package manager brew. 
On Linux, simply use your package manager to install these tools.

\bigskip


To install Orbiter, 
you first need to download the source code from github. 
So, navigate your web browser to
$$
\verb'https://github.com/abetten/orbiter'
$$
Click on the green box that says ``Code'' 
and copy the address into the clipboard 
by clicking the clipboard symbol, see below:
\orbiteroutput{GRAPHICS/LATEX/github}
Click on ``Code'' and copy paste the address. 
Then type 
$$
\verb'git clone XXX' 
$$
Where XXX is to be replaced with the address 
that you copied (simply use CTRL-P to paste).

\bigskip

Once the command finishes, 
you need to check that the prerequisites are satisfied, which are 
\begin{enumerate}
\item
flex
\item
bison
\end{enumerate}
These are compiler generator tools, and they are required for compiling Orbiter. 
The required tools are called bison and flex (in the Unix world, they correspond to yacc and lex, respectively).
Unfortunately, many Mac OS versions come with outdated versions of bison and flex, so 
we need to check the version that is installed 
(or install it in the first place).
To check the version of bison, type 
$$
\verb'bison --version'
$$
If you get a message \verb'bison not found' 
it means that you don't have bison. 
In this case you need to install bison. 
Otherwise, check out the version number. 
Orbiter requires 3.8.2 or higher.
If your version number is below that, 
you need to upgrade your bison software.
Regarding flex, type 
$$
\verb'flex --version'
$$
If you get a message \verb'flex not found' 
it means that you don't have flex. 
In this case you need to install flex. 
Otherwise, check out the version number. 
Orbiter requires 2.6.4 or higher.
If your version number is below that, 
you need to upgrade your flex software.

\bigskip

Next, we will show how to install bison and flex, 
and possibly update the version that is installed.
On a Macintosh system, the package manage brew can be used. 
On Linux, different package managers exists.
In either case, we have to also update certain shell variables, 
so that the correct version of flex and bison 
is used during the Orbiter build process.

\bigskip 

On a Macintosh, using csh (pronounciation: cee-shell), 
you can install bison using ``brew.'' 
Type 
$$
\verb'brew install bison'
$$
The install will go to 
$$
\verb'/opt/homebrew/Cellar'
$$
Go there and check whether a directory bison exists. 
If so, do 
$$
\verb'cd bison'
$$
and 
$$
\verb'ls'
$$ 
and note the version number. 
You may still have an old version of bison somewhere else. 
In order to activate the new bison, you will 
have to change the path in your shell. 
For csh, make sure you add the line
$$
\verb'set path = ( /opt/homebrew/Cellar/bison/3.8.2/bin $path )'
$$
Here, the \verb'3.8.2.' part may has 
to be changed depending on your version of bison. 

\bigskip

You can install flex using brew. Type 
$$
\verb'brew install flex'
$$
The install will go to 
$$
\verb'/opt/homebrew/Cellar'
$$
Go there and check whether a directory flex exists. 
If so, do 
$$
\verb'cd flex'
$$
and 
$$
\verb'ls'
$$ 
and note the version number. 
Like with bison earlier, 
you may still have an old version of flex somewhere else. 
In order to activate the new flex and not the old, 
you will have to change the path in your shell. 
For csh, make sure you add the line
$$
\verb'set path = ( /opt/homebrew/Cellar/flex/2.6.4_2/bin $path )'
$$
Here, the \verb'2.6.4_2' 
part may has to be changed depending on your version of bison. 

\bigskip

Finally, you have to add 
$$
\verb'set LDFLAGS = ( -L/usr/local/Cellar/flex/2.6.4_2/lib -L/usr/local/Cellar/bison/3.8.2/lib )'
$$
to your \verb'~/.cshrc'
Here, the version numbers may have to be 
changed to match those that you used before.

\bigskip

After these changes, 
you need to read the \verb'~/.cshrc' 
file again and rehash. Type
$$
\verb'read ~/.cshrc'
$$
and
$$
\verb'rehash'
$$
If you type 
$$
\verb'bison --version'
$$ 
you should get the correct version number of bison. 
Likewise, if you type 
$$
\verb'flex --version'
$$
you should get the correct version number of flex. 

\bigskip


You are now ready to compile Orbiter. 
For this purpose, you need to enter 
the Orbiter home directory. 
Assuming that you are still 
in the directory where you did \verb'git clone', type 
$$
\verb'cd orbiter'
$$
and then 
$$
\verb'make'
$$
If your computer CPU has multiple cores, you can try 
$$
\verb'make -j 8'
$$
The number 8 can be adjusted 
to the number of cores that you have. 
This multitasking reduces the time it takes to compile Orbiter. 
Once done, you should have an orbiter executable 
in the directory 
$$
\verb'orbiter/src/apps/orbiter'
$$
The executable is called 
$$
\verb'orbiter.out'
$$
Note that you are already in orbiter, so you could do 
$$
\verb'ls src/apps/orbiter'
$$
to check if \verb'orbiter.out' appears.
In either case, you can do 
$$
\verb'cd ..'
$$
to leave the orbiter branch.
You should only go back in 
the Orbiter branch 
if you want to update Orbiter. 
Otherwise, you should not be there. 
In particular, 
you should not change any files in this branch, 
as this would cause problems in the update. 



\bigskip


Finally, if you have to update bison or flex, 
you can type 
$$
\verb'brew upgrade bison'
$$
or 
$$
\verb'brew upgrade flex'
$$


